\documentclass[11pt]{article}

\usepackage[margin=1.1in]{geometry}
\usepackage{amsmath,amssymb,amsthm}
\usepackage{booktabs}
\usepackage{enumitem}
\usepackage{xspace}
\usepackage{xcolor}
\usepackage[colorlinks=true,linkcolor=blue!50!black,citecolor=blue!50!black,urlcolor=blue!50!black]{hyperref}

% ---------------------------------------------------------------------------
% Macros (shared with multi-table-design.tex)
% ---------------------------------------------------------------------------
\newcommand{\F}{\mathbb{F}}
\newcommand{\Ftwo}{\F_2}
\newcommand{\eq}{\mathsf{eq}}
\newcommand{\nnz}{\mathsf{nnz}}
\newcommand{\ip}[2]{\langle #1, #2 \rangle}
\newcommand{\code}[1]{\texttt{#1}}
\newcommand{\out}{\mathrm{out}}
\newcommand{\inn}{\mathrm{in}}
\newcommand{\sid}{s_{\mathrm{id}}}
\newcommand{\ssig}{s_{\sigma}}
\newcommand{\cells}{\mathcal{P}}
\newcommand{\settled}[1]{\medskip\noindent\textbf{SETTLED} (#1).\xspace}

\newtheorem{lemma}{Lemma}
\newtheorem{remark}{Remark}
\theoremstyle{definition}
\newtheorem{definition}{Definition}

\title{Circuits over Flock Tables:\\ Large-Field Gates, Copy-Constraint Wiring, and the Road to Recursion\\[0.6ex]
\large Design document --- working draft v0.2}
\author{Flock}
\date{July 29, 2026}

\begin{document}
\maketitle

\begin{abstract}
This document designs circuit proving over the multi-table core: wires
carry 128-bit words, each gate is one row of a registry table, and the
wiring (plus public inputs/outputs) is enforced on top of the unchanged
per-row relations. The driving goal is \textbf{recursion-completeness}: the
gate inventory must express Flock's own verifier. That goal fixes the
architecture. (i) The registry hosts two \emph{classes} of table in one
union commitment: boolean tables (hashes, GF(2) traces, ring-switched
claims) and \textbf{large-field tables} ($\F_{2^{128}}$-native R1CS, one
constraint per field multiplication, Spartan-style PIOP, packed-direct
claims with no ring-switch) --- a committed word is an $\F$-element either
way, so the classes differ only in their PIOPs, which run on disjoint
regions of the address space. (ii) Wiring is a Plonk-style copy-constraint
permutation realized with the reviewed \code{product\_gkr} grand product;
under the multi-table addressing conventions its evaluation claim
\emph{factors into a handful of packed evaluations of the committed union
polynomial} (Lemma~\ref{lem:gather}) --- no gather sumcheck, no new
committed columns. (iii) Word wires are $\F$-elements under the identity
basis map, which is exactly what makes the hash$\leftrightarrow$field
boundary of the recursion plan a pure copy argument. Status: the
large-field class, its union integration, and the wiring layer are all
\textbf{implemented} --- circuit proofs verify end to end, with the
SHA-256 binary tree as the driving workload.
\end{abstract}

\tableofcontents

% ===========================================================================
\section{Goals and non-goals}\label{sec:goals}

\paragraph{Goal.} One proof attesting that a \emph{circuit} evaluates
correctly: a DAG whose gates are rows of registry tables and whose wires
carry 128-bit words --- (i) every row satisfies its type's relation, (ii)
equal wire endpoints carry equal words, (iii) designated wires equal
statement values.

\paragraph{Goal: recursion-completeness.} The architecture must capture
\emph{the Flock verifier itself} as a circuit. Per the measured taxonomy in
\code{docs/local/recursion-verifier-map.md}, the verifier is
$\sim$10.2k SHA-256 compressions (78--98\% of the circuit: Merkle queries
$\sim$9.1k, FS transcript $\sim$1.0k, PoW just 14 --- grinding is
prover-side work, the verifier checks one hash per grind and the
leading-zero test is free bit logic) plus $\sim$46k
$\F$-multiplications, plus control flow. \S\ref{sec:recursion} traces what this demands; the short
version is that it fixed three choices: the large-field table class
(\S\ref{sec:element}), 128-bit wire granularity (a word \emph{is} an
$\F$-element in the identity basis, so field- and hash-native gates
exchange values by pure copy), and the succinctness-closure requirements
on $\sigma$ and the lincheck comb (\S\ref{sec:sigma},
\S\ref{sec:recursion}).

\paragraph{Design priorities} (inherited from the multi-table document):
prover throughput; gradual landing with byte-identity anchors; static
registry structure with per-proof counts.

\paragraph{Non-goals (v1).}
\begin{itemize}[nosep]
  \item \emph{Dynamic dataflow}: data-dependent addressing, RAM, shared
        lookup tables, prover-chosen wiring (needs LogUp multiplicities;
        \S\ref{sec:logup}). The verifier does \emph{not} need it once the
        inner protocol is shape-normalized (\S\ref{sec:recursion}).
  \item \emph{Sub-word wiring}: quantities finer than 128 bits are not
        routed; sub-word gate parameters are pinned per-type or exposed as
        public constant words; verifier control bits live in boolean-class
        tables where bit access is free.
  \item \emph{Cross-proof buses}: a circuit lives inside one proof.
\end{itemize}

% ===========================================================================
\section{Architecture overview}\label{sec:arch}

Three layers over the unchanged multi-table core:

\begin{enumerate}[nosep]
  \item \textbf{Gate semantics = class-tagged tables}
        (\S\ref{sec:element}; \emph{implemented}). Boolean tables are the
        existing hash tables. Large-field tables carry $\F$-native R1CS
        rows (multiplication gates and friends). One union commitment
        hosts both; the two PIOP families run on disjoint address regions
        and meet at the transport batch.
  \item \textbf{Wiring = copy constraints} (\S\ref{sec:wiring};
        \emph{implemented}). A permutation $\sigma$ on the cell space of all
        IO words, checked by \code{product\_gkr}; the wire-value claim
        factors into packed-direct claims on the union polynomial.
  \item \textbf{Public IO = a public cell segment} (\S\ref{sec:cells}):
        statement words and constants are cells whose values the verifier
        evaluates itself; $\sigma$ wires them into the gates. Constants
        (IVs, the constant-1 cell that pins inversion gadgets) come for
        free.
\end{enumerate}

% ===========================================================================
\section{Preliminaries}\label{sec:prelim}

We use the multi-table document's notation: uniform row capacity $2^\nu$,
union bit-address space of $M$ variables, packed polynomial $\hat f$ over
$M-7$ word variables, committed dense stack $\hat q$ (height-$n_t$
stacking; dropped words are definitionally zero and the jagged transport
enforces it). Facts this design leans on:

\begin{enumerate}[nosep]
  \item \textbf{BatchMajor, rows low.} The packed word address of
        (chunk-column $w$, row $i$) inside a slot is $(w \ll \nu) + i$.
        One committed $\F$-element is 128 consecutive trace bits of a
        single invocation (boolean class) or one witness element
        (large-field class).
  \item \textbf{Slot alignment.} A packed address decomposes as
        $[\, i \mid w \mid p_t \,]$ with the slot prefix $p_t$ Boolean.
  \item \textbf{128-aligned IO regions.} Every wireable region of every
        encoder sits on a word boundary.
  \item \textbf{Claim seam.} Throughout, \emph{the transport} means the
        multi-table document's \textbf{jagged transport} --- the
        commitment-side opening pipeline ($\gamma$-batching $\to$
        virtual-open/merged sumcheck $\to$ the jagged identity
        $\ip{q}{W_\rho} = \hat f(\rho)$ $\to$ one Ligerito opening) that
        \emph{consumes} every PIOP's claims; it is not the wiring/bus
        argument, which is a claim \emph{producer} upstream of it
        (\S\ref{sec:wiring}). The transport is class-agnostic: it
        $\gamma$-batches \emph{weighted openings of the committed words}.
        Each table class prepares its own claims. Bit structure is
        internal to the boolean class --- its constraints are over bits,
        so its PIOP claims exit through ring-switching, \emph{that
        class's} bit$\to$word bridge --- while the element class and the
        wiring layer emit word-level evaluations directly
        (\code{PackedDirectClaim}s, now carried on the union path). Bits
        never cross the seam; what crosses is a (weight, value) pair. The
        one trace of bit-ness the transport carries is that the boolean
        bridge's weights are Frobenius-structured (128 components) rather
        than plain $\eq$ --- the element/wiring weights are the
        degenerate single-component case (\S\ref{sec:routing}).
\end{enumerate}

\subsection{The product-GKR primitive (reviewed)}\label{sec:gkr}

\code{product\_gkr} proves, for $f, g : \{0,1\}^\mu \to \F$ and
$\sigma : [N] \to [N]$, $N = 2^\mu$, with challenges $\alpha, \beta$
squeezed at entry:
\[
  \prod_i \bigl(f_i + \alpha\,\sid(i) + \beta\bigr) \;=\;
  \prod_i \bigl(g_i + \alpha\,\ssig(i) + \beta\bigr),
  \qquad \ssig(i) = \sid(\sigma(i)),
\]
certifying the multiset equality $\{(f_i, i)\} = \{(g_i, \sigma(i))\}$;
with $f = g = w$ this is the copy-constraint statement (\emph{$w$ constant
on the cycles of $\sigma$}). \code{prove\_batched} runs both product trees
in lockstep ($\lambda$-combined, eq-weighted degree-2 layers) and lands at
a \emph{single} point $\rho$, emitting $\hat f(\rho), \hat g(\rho),
\hat s_\sigma(\rho)$. No committed oracle, no inversions; prover $O(N)$
multiplications; transcript $O(\mu^2)$; measured 7.6\,ms / 7.4\,KiB at
$\mu = 20$.

The module is a PIOP: its claims are sound only once the caller binds them
--- the statement into the challenger before the call, and the surfaced
evaluations to the committed witness afterwards. The specific requirements
appear where the wiring protocol uses them (\S\ref{sec:wiring},
\S\ref{sec:statement}).

% ===========================================================================
\section{The large-field table class}\label{sec:element}

\emph{Status: implemented} — standalone (\code{5339bee}, \code{bc6bab0})
and integrated into the union (\code{c87a067..09918ed}); mixed
boolean+element proofs verify end to end.

\settled{Ron} Each registry type is class-tagged \code{boolean} or
\code{large-field}. The two PIOPs operate on \textbf{disjoint regions} of
the trace: class-major slot sort, the boolean region a prefix subcube
$[0, 2^{M_{\mathrm{bool}}})$, the element region a disjoint aligned
subcube with Boolean-frozen prefix coordinates. (The one-domain
alternative is \emph{unsound} under the boolean $C{=}I$ convention: the
zerocheck's $c$-side aliases the witness, so nonzero element data inside
the boolean domain breaks the honest sum. The implementation runs the
boolean PIOP over its subcube and appends frozen-zero high coordinates to
its claims.)

\begin{remark}[Dead addresses: the audit, and the plan]\label{rem:dead}
The two-subcube structure inflates the address space: since neither
class region can live inside the other's subcube (the $c{=}z$ trap
again), \emph{any} nonempty element class costs one extra address bit
over the boolean extent --- up to $2\times$ address space, mostly dead,
regardless of how small the element tables are (and symmetrically if the
classes swap order). The audit of who pays for a dead address:
commitment, \emph{nothing} (height-$0$ columns commit no words ---
measured exactly flat); both PIOPs, \emph{nothing} (the dead region is
outside both domains --- the point of the design); wiring,
\emph{nothing} (the cell space indexes (table, io-slot) $\times$ rows
compactly and never sees union addresses); verifier, a few frozen
coordinates. Two places do pay: the \emph{unmerged} transport's
padded-domain auxiliaries ($b_{\mathrm{combined}}$ and its fold chain
scale with $2^{M-7}$ --- a few ms per extra bit at $M \approx 30$), and
the materialized padded witness buffer ($\sim$134\,MB of zero pages per
extra bit at that scale, plus its fill pass). The plan: accept now;
the \textbf{merged transport computes over the dense domain} and has
neither cost, so merged intake (\S\ref{sec:status}) retires the first
payer for circuit proofs; the buffer cost is removable by the
previously-shelved dense-born witness refactor, whose value grows with
exactly this inflation --- revivable if profiles say so. Rejected:
packing the boolean domain as multiple subcubes --- buys back at most
one bit for extra zerocheck transcripts or a non-factoring domain
indicator.
\end{remark}

\begin{remark}[Taxonomy of emptiness]\label{rem:empty}
Four kinds, two index spaces, one invariant. (1) \emph{Non-IO trace
columns}: real committed witness, just not wireable --- no cells index
them. (2) \emph{Dead addresses} (unused block columns, alignment gaps,
the inter-class region): height 0, never committed, outside both PIOP
domains, not cells; the virtual $\hat f$ is definitionally zero there
and the transport identity $\ip{q}{W_\rho} = \hat f(\rho)$ is the
binding. (3) \emph{Dummy rows} ($n_t \le j < 2^\nu$): uncommitted and
virtually zero on the union side (a claim implying a nonzero one breaks
the transport identity --- the closure tests); inside the PIOP domains,
where zero rows satisfy both classes' relations (homogeneity;
$a_0 \odot b_0 = 0$) and run lists skip the work; and dummy \emph{cells}
on the wiring side --- value 0 read from the padded buffer,
$\sigma$-fixed by validation, identical product factors both sides
(the capacity tax is their only cost), zero contributions to the gather
folds. (4) \emph{Cell-space padding} (pad slots, unused public rows):
indices with no address at all --- value 0 by definition, $\sigma$-fixed,
nothing on disk to contradict. The invariant: everything unnamed is
zero, and zero is both \emph{satisfying} (for the relations) and
\emph{inert} (for the wiring) --- enforced everywhere a prover could
cheat by the commitment simply not containing those words.
\end{remark}

\subsection{The relation}

A large-field type is a base block over $\F$: $2^\kappa$ witness columns
per row, \textbf{C = identity} (constraint domain $\equiv$ column domain),
sparse $A_0, B_0 \in \F^{2^\kappa \times 2^\kappa}$, and \textbf{affine
constant vectors} $a_0, b_0 \in \F^{2^\kappa}$ carried in the registry
(not the witness):
\[
  (A_0 z + a_0)\,(B_0 z + b_0) \;=\; z
  \quad\text{per row, per column.}
\]
\settled{Ron} No constant-1 column (it would waste a word-column in
tables only a few columns wide); constants live in the matrices, and the
verifier evaluates $\hat a_0, \hat b_0$ in closed form ---
count-independent, since the constant vector is uniform across rows.
\textbf{Validity rule}, enforced at construction: $a_0 \odot b_0 = 0$
(disjoint supports), so an all-zero row satisfies the relation ---
dummy/padding rows stay complete, sound, and free. Standard rows:
multiplication ($A_0 = e_a$, $B_0 = e_b$), free wire ($A_0 = e_y$,
$b_0[y] = 1$: a tautology), linear-pinned-to-wire ($A_0 =$ combination,
$b_0[y] = 1$), padding column (all-zero row forces $z_y = 0$). Inversion:
constrain $u = x \cdot y$ and wire $u$ to the public constant-1 cell.
Gate granularity beyond \code{mult}/\code{mult-acc} is deferred.

\subsection{The PIOP: Spartan-style, two phases}

\settled{Ron} Phase 1: a large-field zerocheck,
$\sum_x \eq(\tau,x)\bigl((A z + a_0)(Bz + b_0) - z\bigr)(x) = 0$ over the
element region --- plain degree-3 eq-weighted rounds, no univariate skip,
no packing; one zerocheck across all element slots (the same eq-derived
implicit batching as the boolean union). Because $C = I$, the $c$-claim is
\emph{directly} a witness evaluation at the zerocheck's point. Phase 2: an
$\alpha$-batched lincheck for the $A$/$B$ claims.

\paragraph{The collapsed lincheck.} Because the full system is
$I_{2^\nu} \otimes A_0$, the lincheck weight factors as
$\eq(r_{\mathrm{row}}, y_{\mathrm{row}})\cdot \mathrm{comb}(y_{\mathrm{col}})$,
so the row half sums away in closed form: the prover partial-folds $z$ at
$r_{\mathrm{row}}$ (one pass) and sumchecks only the $\kappa$ column
variables. Output claim $\hat z(r_{\mathrm{row}}, r'_{\mathrm{col}})$ ---
row coordinates \emph{inherited} from the zerocheck, so the element
class's two claims share their row point; the verifier's final check is
$(\hat A_0 + \alpha \hat B_0)(r_{\mathrm{con}}, r'_{\mathrm{col}})$ in
$O(\nnz)$ --- a few entries per gate type, so \textbf{the comb problem
does not exist for element tables}. Anything built on top should assume
$\kappa$ rounds. (Landed as \code{c87a067}.)

\paragraph{Witness and padding.} $Az, Bz$ are sparse gathers emitted
during generation (for a mult gate, literally the operand values); dummy
rows are all-zero, satisfy the relation, contribute nothing to round
polynomials, and --- under the union's height-$n_t$ stacking --- are
\emph{not committed}, making dummy-row zero-ness structural (the
$\ip{q}{W_\rho} = \hat f(\rho)$ transport identity breaks otherwise;
test-pinned in both debug and release forms).

\subsection{Claim routing through the transport}\label{sec:routing}

Element claims are packed evaluations --- \emph{no ring-switch}. On the
unmerged jagged path they ride as \code{PackedDirectClaim}s (implemented:
four claims --- boolean AB, C ring-switched; element C, LC packed-direct
--- through one opening). On the \textbf{merged} (v6) transport the same
composition holds structurally: the merged sumcheck is a $\gamma$-batched
multi-weight transport whose per-claim weights are
\code{FrobeniusClaim}\{$z_{\mathrm{row}}, z_{\mathrm{col}}$,
coeffs[128]\}; a ring-switched claim contributes its 128-coefficient
Frobenius residue, and an element claim is the \emph{degenerate}
single-component case $\mathrm{coeffs} = \gamma \cdot e_0$ --- no extra
rounds, weight factors as $\eq_{\mathrm{row}} \otimes \eq_{\mathrm{col}}$
(streamable). A pure-element proof degenerates to plain
jagged$\to$Ligerito with no Frobenius content at all. \emph{Extending the merged path's claim intake --- today it accepts
ring-switched claims only --- is future work}; known items: the point
split assumes a skip coordinate, the assist should specialize
single-component claims (ties into the \emph{jagged remainder},
\S\ref{sec:open}),
\code{stream\_b} currently disables on packed-direct claims (accepted,
TODO), and the merged entries assert element-free until then. Until that
lands, element and circuit proofs ride the unmerged path.

% ===========================================================================
\section{The circuit model}\label{sec:cells}

\begin{definition}[IO schema]
A registry type carries lists of 128-aligned trace-word positions
designating gate inputs and outputs (generalizing \code{ChainLayout});
words outside the schema are internal. \emph{Implemented}
(\code{schedule::IoWord}); the payload appends to the registry digest only
when non-empty, so schemaless registries digest byte-identically.
\end{definition}

\begin{definition}[Circuit]
Gate counts per type (each gate one row), a set of \emph{public cells}
(statement words and constants), and a wiring: a partition of all IO cells
into wires. Valid iff every wire has at most one producer, the dataflow
graph is acyclic, and the induced $\sigma$ is a permutation.
\end{definition}

\begin{remark}[Semantics are satisfiability]
The proof enforces gates + wiring; that a satisfying assignment computes
``the'' output is a property of the circuit (single producer, acyclicity),
checked by whoever endorses the circuit digest — as in Plonk.
\end{remark}

\paragraph{The cell space.} Enumerate cell-slots: pairs (type, io-word
position) across both classes, plus public cell-slots, padded to $2^c$.
$\cells = \{0,1\}^{\mu}$, $\mu = \nu + c$, cell $p = (j, \iota)$ with the
row $j$ in the \emph{low} bits. A gate cell maps to committed word
$W_\iota(j) = [\, j \mid \mathrm{chunk}(\iota) \mid p_{t(\iota)}\,]$;
public slots hold statement words; everything else is a dummy cell (value
0, fixed by $\sigma$). Wires are $\sigma$-cycles: fan-out $k$ is a cycle
of length $k{+}1$; a constant wired into $n$ gates is one long cycle;
unread outputs are fixed points. The wire-value polynomial $w$ assigns
each cell its word.

Worked example: $\nu = 10$, a SHA-256 gate table with schema words
$(h_0, h_1, m_0 \ldots m_3, o_0, o_1)$, a mult table with $(a, b, c)$,
one public slot --- $12$ cell-slots padded to $2^c = 16$, so $\mu = 14$
and $16{,}384$ cells. (The schema is a per-type choice, not forced by
the trace: this is the \emph{chaining} compression gate, $h$ wireable so
multi-compression hashes --- FS chains, Merkle leaves --- wire gate
$i{+}1$'s $h$ to gate $i$'s output, and IV instances wire $h$ to the
public IV cell. An IV-\emph{pinned} 6-word variant drops $h_0, h_1$ and
pins the IV in the relation instead; a word left out of the schema must
be pinned by the relation, else it is free witness.) Cell $(\iota{=}2, j{=}17)$ is ``the $m_0$ input
word of sha2 gate \#17''; its committed word is the bit-concatenation of
$\iota$'s two registry constants with $j$:
$[\,p_{\mathrm{sha2}} \mid \mathrm{chunk}(m_0) \mid j\,]$. Everything in
the wiring argument speaks this $14$-bit language --- $w$, $\sigma$, the
$\sid$ tags in the product are \emph{cell} indices, never union
addresses. The point of the second index space is compression: run over
union addresses (with $\sigma$ fixing non-IO words) and the product pays
a factor per \emph{trace word} --- the sha2 slot alone is $2^{18}$
words, dominated by round intermediates no wire touches, plus the dead
inter-class space of Remark~\ref{rem:dead}; the cell space pays per
\emph{schema word}, $2^{14}$ here, and is immune to both. The verifier's
per-slot data is tiny: $2^c$ coefficients $\eq(\rho_\iota, \iota)$ and
each $\iota$'s constant (prefix, column) pair from the registry.

% ===========================================================================
\section{The wiring argument}\label{sec:wiring}

\subsection{Mechanism: why a permutation proves the copy
constraints}\label{sec:wiring-mech}

Terminology first: ``copy constraints'' is Plonk's name for the
\emph{wiring equalities} --- the requirements that certain cells hold
equal values. They are \textbf{not} R1CS constraints: no gate, row,
column, or matrix entry implements them, and the R1CS stays
block-diagonal with each gate an island. Their entire enforcement is the
product identity below, which lives outside the constraint system and
commits nothing.

\paragraph{No dataflow.} At the proof level nothing is ``read'' or
``routed''. Every gate is a self-contained row whose constraints touch
only its \emph{own} cells. ``Gate $A$'s output feeds gates $B$ and $C$''
is purely the statement that three committed cells --- $A$'s output cell
$p$, $B$'s input cell $q$, $C$'s input cell $r$ --- \emph{hold equal
values}. A wire is an equivalence class of cells; fan-out $k$ is a class
of size $k{+}1$. The value is not moved anywhere; it is present in
several cells, and the wiring argument proves those cells agree.

\paragraph{Classes as a permutation.} Pick any permutation $\sigma$ of
$\cells$ whose orbits are the wire classes --- concretely, rotate each
class cyclically ($\sigma(p) = q,\ \sigma(q) = r,\ \sigma(r) = p$ in the
example) and fix every unwired cell.

\begin{lemma}[Tag rigidity]\label{lem:tags}
Let $L = \{(w_x, x)\}_{x\in\cells}$ and
$R = \{(w_x, \sigma(x))\}_{x\in\cells}$, as multisets of
(value, position-tag) pairs. Then $L = R$ iff $w$ is constant on every
orbit of $\sigma$.
\end{lemma}

\begin{proof}
Each side holds exactly one pair carrying each tag, so multiset equality
must match them tag-by-tag: the $L$-pair at tag $q$ is $(w_q, q)$, while
the $R$-pair at tag $q$ is $(w_{\sigma^{-1}(q)}, q)$. Hence $L = R$ iff
$w_q = w_{\sigma^{-1}(q)}$ for every $q$, i.e.\ $w$ is
$\sigma$-invariant, i.e.\ constant on orbits. In the example: tag $q$
forces $w_q = w_p$, tag $r$ forces $w_r = w_q$, tag $p$ forces
$w_p = w_r$ --- chasing tags around the cycle equates the whole class,
which is exactly how one cycle of length $k{+}1$ enforces a fan-out-$k$
wire.
\end{proof}

\paragraph{Multisets as products.} Encode a pair as the linear factor
$(w_x + \alpha\,\sid(x) + \beta)$ in formal variables $\alpha, \beta$.
The map pair $\mapsto$ factor is injective and polynomials factor
uniquely, so
\[
  \prod_x \bigl(w_x + \alpha\,\sid(x) + \beta\bigr)
  \;=\;
  \prod_x \bigl(w_x + \alpha\,\sid(\sigma(x)) + \beta\bigr)
\]
holds \emph{as polynomials in $\alpha,\beta$} iff $L = R$. The verifier
samples $\alpha, \beta$ \textbf{after} the witness is committed and
checks the two products at that one point; a false identity survives with
probability $O(N/2^{128})$ (Schwartz--Zippel). Note both sides use the
\emph{same} committed value vector $w$ --- only the tag pairing differs.
This is the $f{=}g{=}w$ instantiation of \S\ref{sec:gkr}, and computing
and comparing the two grand products is precisely what the GKR does,
with no committed oracle.

Corollaries: direction never enters --- $\sigma$ neither knows nor cares
which cell is the producer; an output read by nobody is a fixed point
(no constraint, no cost); single-producer and acyclicity are
circuit-\emph{validation} properties (they make the satisfying
assignment unique given the inputs), not proof obligations.

\subsection{The protocol and the gather factorization}

After commitment and statement binding, run
\code{product\_gkr::prove\_batched}$(f{=}g{=}w, \sigma)$ over $\cells$:
all wiring and public-IO equalities at once, claims $\hat w(\rho)$ and
$\hat s_\sigma(\rho)$ at one $\rho = (\rho_r, \rho_\iota)$. Dummy cells
cost nothing and threaten nothing ($\sigma$-fixed, identical factors both
sides).

\begin{lemma}[Gather factorization]\label{lem:gather}
With cells indexed rows-low,
\[
  \hat w(\rho_r, \rho_\iota)
  \;=\;
  \sum_{\iota\ \mathrm{gate}} \eq(\rho_\iota, \iota)\cdot
     \hat f\bigl(\rho_r \| \mathrm{chunk}(\iota) \| p_{t(\iota)}\bigr)
  \;+\;
  \sum_{\iota\ \mathrm{public}} \eq(\rho_\iota, \iota)\cdot
     \hat v_\iota(\rho_r).
\]
\end{lemma}

\begin{proof}
Expand and use multilinearity of $\hat f$ in its low (row) coordinates:
$\sum_j \eq(\rho_r, j) \hat f(j \| x) = \hat f(\rho_r \| x)$. Public slots
contribute their MLEs; dummies contribute 0.
\end{proof}

Every gate term is a packed evaluation at a point random in the row
coordinates and Boolean in all high coordinates --- the existing
packed-direct claim shape, consumable with no gather sumcheck and no new
committed columns. The $\eq(\rho_\iota, \iota)$ coefficients and the
public sum are verifier-computed in $O(2^c + \#\mathrm{public})$; the
verifier checks the transported values recombine to $\hat w(\rho)$,
\textbf{and} that
$f_{\mathrm{eval}} = g_{\mathrm{eval}} = $ this recombined value ---
\emph{both} of the GKR's surfaced evaluations must be bound to the
committed witness; binding only one leaves the other product's input
vector prover-chosen. Likewise, v1 verification uses the $\sigma$-aware
verifier (never the trusting variant), and circuit validation enforces
that $\sigma$ is a permutation (completeness; soundness fails closed).

\begin{remark}[Where each convention is load-bearing]
Rows-low lets $\hat f$'s multilinearity absorb the row fold; uniform
capacity makes a slot's row variable the union row variable; slot
alignment freezes the high coordinates Boolean; 128-aligned IO regions
make a wire word one committed word. Remove any one and the lemma degrades
into a real gather argument.
\end{remark}

\paragraph{Capacity tax.} \settled{Ron; recorded as potential
improvement} The product runs over all $2^\mu$ cells --- capacity, not
counts: a dummy factor $\alpha\sid(p) + \beta \ne 1$ cannot be skipped. At
$\mu \approx \nu + 5$ this is $\sim$12\,ms at $\nu = 16$; sizing $\nu$ to
the circuit family makes it the honest cost. The v2 refinement (mask dummy
leaves to 1 --- sound since $\sigma$ maps live cells to live cells; all-1
subtrees skip by run lists; count-derived closed-form corrections on the
verifier side) is sketched and deferred. The 1's live only in the
prover's uncommitted product-tree leaf vector, built per proof from the
committed values and $\alpha, \beta$ --- the trace's dummy rows stay
all-zero, as the zerocheck and jagged transport require (indeed the
masked input-layer reconstruction uses $\hat w_{\mathrm{live}} = \hat w$,
which holds \emph{because} dummy cells read zeros); the correction terms
are the $\mathrm{live}$-indicator MLEs, per-slot aligned-interval eq
sums the verifier evaluates from the counts.

% ===========================================================================
\section{Binding \texorpdfstring{$\hat s_\sigma(\rho)$}{s-sigma}: the
circuit lifecycle}\label{sec:sigma}

\settled{Ron} \textbf{v1: verifier-known circuit.} The verifier holds the
circuit and evaluates $\hat s_\sigma(\rho)$ in $O(|\cells|)$
(\code{verify\_batched\_with\_sigma}) --- a few ms, subdominant to the
lincheck comb the verifier already pays ($O(\nnz)$, $\sim$83\% of verify).
\textbf{v2: preprocessed $\sigma$ commitment} for repeated circuits ---
the \emph{same} delegation problem as comb delegation (registry-/circuit-
static data, commit once, open per proof); the two must share one design.
Recursion \emph{requires} the v2 form (an $O(2^\mu)$ in-circuit evaluation
is unaffordable); v1 is for native verification. Structured wirings
(chains, complete trees) evaluate in $O(\mu)$ --- the existing
\code{chain.rs}/\code{merkle\_path.rs} glue is the degenerate case of this
layer and stays.

% ===========================================================================
\section{Why copy constraints and not a LogUp bus}\label{sec:logup}

For a \emph{fixed} circuit, $\sigma$-cycles handle fan-out natively with a
pure permutation and zero new committed columns, using the reviewed
primitive. A LogUp bus (emit/consume tuples with multiplicities) is the
right architecture for \emph{dynamic} dataflow --- and \code{logup\_gkr}
exists unported on \code{flock-dev} --- but the verifier does not need it
(\S\ref{sec:recursion}), and the migration path preserves everything: the
cell space, IO schemas, Lemma~\ref{lem:gather} and the transport seams
carry over; only the multiset argument swaps. Prover-\emph{ordered} data
(sorted query lists) needs only plain multiset equality --- no tags, no
trust boundary --- which the same primitive provides.

% ===========================================================================
\section{Recursion: capturing the verifier}\label{sec:recursion}

The verifier is expressible as a \emph{fixed-topology} circuit in this
model, given three things (all now designed or landed):

\begin{enumerate}[nosep]
  \item \textbf{The large-field class} (\S\ref{sec:element};
        \emph{landed}): the $\sim$46k multiplications cost 1 constraint
        each instead of $\sim$2187 bit-arithmetized ones. The recursion
        map's ``two instances composed by a crossing argument'' becomes
        two \emph{regions of one committed trace}, and the crossing
        ($\sim$15.8k opened-row elements + transcript) is ordinary wiring
        --- pure copy, because the bit$\leftrightarrow$element map is the
        identity in the standard basis. Element slots need no
        ring-switch, so the verifier-of-the-verifier also loses the
        $\code{eval\_rs\_eq}$ family ($\sim$15k mults): recursion gets
        simpler as it recurses.
  \item \textbf{Inner-protocol shape normalization} (family J), settled:
        \settled{Ron} query sampling \emph{with replacement} ---
        \emph{landed}, see \S\ref{sec:replacement-sampling}; the rejection
        loop, the \code{HashSet} and the sort are gone, and the soundness
        re-accounting came out at \emph{no change to any query count}.
        Merkle openings as \emph{full independent fixed-depth paths}
        first ($\sim$30\% more hash work, fully static circuit), with
        \emph{Merkle caps} (commit the top-$2^k$ layer; every path
        shortens by $k$) as the later knob --- still owed. \settled{Ron}
        The Merkle/FS hash is \textbf{BLAKE3 throughout} --- the
        $2\times$ circuit-size win inverts the native HW-SHA decision;
        the selectable-hash machinery is merged.
  \item \textbf{Succinctness closure}: every verifier component polylog
        in-circuit or delegated.
        \textbf{The matrix work --- both classes' combs --- is SOLVED, and
        by deferral rather than delegation}: the verifier emits its
        bilinear-form claims instead of evaluating them, and they fold into
        a constant-size accumulator discharged once at the root. That is a
        design in its own right; see
        \code{docs/folding-verifier-design.tex}, which also carries the
        argument for why arithmetising the comb directly can never close a
        fixed point ($\nnz$-preserving), the merge composition, and the
        retraction that this buys nothing natively. Landed, with a native
        merge node folding two circuit proofs.
        Remaining violations: the v1 $\sigma$ evaluation
        ($O(2^\mu)$ --- v2 fixes it) and the \emph{jagged remainder}
        (\S\ref{sec:open}), which is no longer avoidable by choosing a
        transport now that circuit proofs ride the merged one. The wiring
        layer's own verifier is $O(\mu^2)$ sumcheck replay --- closure-safe.
\end{enumerate}

Bit-level control (query index bits, path direction selects, FS byte
packing) lives in boolean-class glue tables where $\Ftwo$-linear work is
free; word granularity carries everything across the boundary.

\subsection{Query sampling with replacement}\label{sec:replacement-sampling}

The inner protocol used to sample \emph{distinct} query positions:
draw, reduce mod the block length, reject on a repeat, and finally sort.
An unbounded loop, a hash set and a sort --- three shapes a
fixed-topology circuit cannot express. The sampler now draws with
replacement: exactly $q$ draws, duplicates kept, returned in sample
order. In circuit terms it is one fixed-length squeeze followed by a
low-bit mask per element, and nothing else. (Block lengths are powers of
two, so the mask is exact --- the reduction was never a source of bias,
only the distinctness was a source of control flow.)

\paragraph{The re-accounting: no query count moves.}
Let $S$ be the agreement set of a word at distance $\gamma$ from the
code, $|S|/n \le 1-\gamma$. The bad event is that every queried position
lands in $S$. For $q$ independent uniform draws that probability is
exactly $(|S|/n)^q \le (1-\gamma)^q$. Sampling without replacement gives
the hypergeometric $\prod_{i<q} (|S|-i)/(n-i)$, which is
\emph{smaller}. The implementation solves $q = \lceil 100 / \log_2
(1/(1-\gamma)) \rceil$ --- i.e.\ it was already pricing the
with-replacement bound, and the old distinct-query sampler was buying
slack that no formula ever spent. So the switch costs nothing and every
per-level query count is unchanged: 243 at rate $1/2$, 148 at $1/4$, and
so on. The $\alpha$-batching of the induced basis is likewise
indifferent --- a repeated position merges into one constraint carrying
the summed weight, and a violated position still puts a nonzero entry at
both slots, leaving the multilinear Schwartz--Zippel term $\lceil \log_2
Q\rceil / 2^{128}$ untouched. The one real cost of a duplicate is one
fewer independent position, and that is exactly what the bound above
already charges for.

\paragraph{Where the obligation moved.}
Two, both native-side. The octopus multi-proof wants strictly ascending
unique leaf positions, so the verifier now collapses the multiset itself
--- and while collapsing it must \emph{reject a prover that answered one
position with two different rows}, since only one row per position
reaches the Merkle check while every slot feeds the induced basis. That
sort does not need to arithmetise: it disappears in-circuit under the
other half of family J, where openings become independent fixed-depth
paths and a duplicate is simply a repeated path.

\paragraph{Consequence for the FS chain.}
The draws are now a single batched squeeze per level rather than $q$
sequentially-dependent duplex rounds. Measured over a whole verification
(\code{--features hash-count}, \code{verifier\_hash\_count}), with the
BLAKE3 estimate split into its four parts:

\begin{center}
\begin{tabular}{lrrrrr}
\textbf{m=30, rate 2} & absorb & parent & \textbf{finalizations} & xof out & total \\
\hline
distinct sampling      & 295 & 19 & \textbf{363} & 123 & 800 \\
with replacement       & 287 & 18 & \textbf{106} & 123 & 534 \\
\end{tabular}
\end{center}

\noindent
The entire win is in the \emph{finalizations}, and that is the right
place for it: absorb blocks are sequential only within a 1\,KiB chunk,
XOF output blocks are counter-mode and mutually independent, but each
squeeze's output is re-absorbed --- so the finalization count \emph{is}
the serial depth of the transcript. At L0 specifically, 243
finalizations become one, and the 243 XOF output blocks become 61,
because the old path squeezed 16 bytes at a time and discarded 48 of
every 64. Sizing the FS chain table off the total (534) rather than the
parts would miss all of this.

A bit-packed squeeze (one field element carrying $\sim$6 query indices)
would cut the XOF-output and absorb terms further and remains available;
it was not taken because it trades the short low-bit mask for a full
128-bit decomposition per squeezed element, and because those terms are
no longer the binding ones.

The ladder shapes still refuse a block narrower than its query count.
That is now a proof-size convention rather than a soundness requirement
--- such a block would open most of itself --- and holding it fixed is
what keeps every shipped ladder identical across the change.

% ===========================================================================
\section{Statement, transcript, Fiat--Shamir}\label{sec:statement}

Statement: (registry digest, circuit digest, public words, commitment
root); the circuit digest covers gate counts (which determine the union
counts), IO schemas, wiring, and the public-cell layout. Transcript order:
commit $\to$ \code{bind\_statement} $\to$ boolean $\tau$ $\to$ boolean
zerocheck $\to$ boolean lincheck $\to$ element $\tau'$ $\to$ element
zerocheck $\to$ $\alpha'$ $\to$ element lincheck $\to$ wiring GKR
($\alpha, \beta$ are squeezed at its entry, so statement binding must
precede it) $\to$ $\gamma$-batched
opening. \emph{Implemented end to end}: the circuit half of the binding
(digest + public words) appends under its own versioned label, so a
non-circuit proof's transcript is unchanged byte for byte.

% ===========================================================================
\section{Cost accounting}\label{sec:costs}

Element PIOP: sub-ms at verifier-circuit scales; claims add $O(1)$ terms
to the combine. Wiring: $O(2^{\mu})$ prover multiplications
($\mu = \nu + c$, $\sim$12\,ms at $\mu = 21$), transcript +7--15\,KB;
verifier $O(\mu^2)$ + $O(2^c + \#\mathrm{public})$ + (v1 only)
$O(|\cells|)$ for $\hat s_\sigma$. Measured baselines and the recursion
sizing ($m \approx 26.6$ with BLAKE3 + Slim ladder) live in the recursion
map and the multi-table document.

% ===========================================================================
\section{Implementation status and roadmap}\label{sec:status}

\paragraph{Landed (2026-07-29, branch \code{multitable}).}
\begin{itemize}[nosep]
  \item BLAKE3 IO word alignment (\code{f95dfbb}) --- fixtures re-pinned;
        SHA-256 anchor byte-identical across the change.
  \item Large-field table class, standalone (\code{5339bee},
        \code{bc6bab0}): types, both PIOP phases, end-to-end single-table
        proof, tamper matrix.
  \item Union integration (\code{c87a067..09918ed}): collapsed lincheck,
        class tags (boolean-only digests byte-identical), class-major
        layout, disjoint PIOPs, four claims through one unmerged opening,
        dummy-row closure. All existing anchors green untouched.
  \item \code{product\_gkr} soundness review (two passes, 2026-07-29).
  \item \textbf{Wiring milestone} (2026-07-30, \code{6b1dc6a..f1c8721}):
        registry IO schemas, the circuit module (cell space,
        validation, canonical wiring, digest, $\sigma$), the wiring GKR
        and gather recombination, public IO, and the circuit prove/verify
        entries --- gather claims joining the element claims' unmerged
        packed-direct intake, no transport change. All four binding
        obligations of \S\ref{sec:wiring}/\S\ref{sec:statement} are
        test-pinned, the $g$-side one by running the real $g = w \circ \sigma$
        forgery. Driving workloads green: a SHA-256 binary tree against a
        native tree computation, an element chain, the cross-class
        crossing, a randomized-wiring differential oracle. Every existing
        anchor and byte fixture untouched. Measured at $\nu = 10$ over 511
        SHA-256 gates ($\mu = 14$): wiring GKR + gather $1.2$\,ms of a
        $12.1$\,ms prove, $4.1$\,KB of a $294$\,KB proof, 8 gather claims.
\end{itemize}

\paragraph{Next.}
\begin{enumerate}[nosep]
  \item \textbf{Delegation design} (comb + $\sigma$, one design) --- the
        recursion map's top-priority item.
  \item \textbf{Merged-transport intake} for packed-direct claims +
        \code{stream\_b} support (perf follow-up; unmerged path is the
        correctness oracle).
  \item Deferred: count-proportional element prover; LogUp bus;
        capacity-tax refinement; Merkle caps; keccak batch-major tier.
\end{enumerate}

% ===========================================================================
\section{Open questions}\label{sec:open}

\begin{enumerate}[nosep]
  \item Delegation design shape (comb + $\sigma$ + assist): SPARK-style
        sparse commitment vs.\ alternatives; adaptive-ordering caveat from
        the recursion map (\S5.1) must be preserved.
  \item Element gate granularity: \emph{candidate list derived}
        2026-07-31 (see \code{docs/local/recursion-handoff.md},
        ``The macro-gate operation list''). Scoping the replay path
        ($\sim$1{,}370 lines, once deferral removes the matrix work) yields
        $\sim$19 operations in three tiers: scalars
        (\code{add}/\code{mul}/\code{inv}/\code{assert\_eq}), sumcheck-round
        machinery, and tensor ops (\code{eq\_*}, Lagrange/interpolation, the
        induced-basis evaluations). The same list is the \code{Arith} trait
        the verifier would be written against, so choosing gate granularity
        and choosing the abstraction are one decision.

        Granularity should be \emph{macro}, not add/mul, and there is now a
        concrete reason rather than an aesthetic one. $\varphi_8$ is a field
        homomorphism, so the interpolation node set
        $V = \varphi_8(\{0..2^m-1\})$ is an $\Ftwo$-subspace and its vanishing
        polynomial $Z_V(X) = \prod_{s \in V}(X+s)$ is \emph{linearized}
        (verified: \code{phi8\_node\_sets\_have\_linearized\_vanishing\_polynomials}).
        Hence $L_i(z) = Z_V(z) / ((z + \varphi_8(i))\cdot d_i)$ with $d_i$
        constant --- one additive map and one inverse in place of an
        $O(\ell^2)$ product. At $\ell = 64$ the combined interpolation's
        $\sim$16k multiplications become $\sim$128, and $Z_V$ being
        $\Ftwo$-linear makes it \emph{free in the boolean class}. Natively the
        naive product is sub-millisecond and not worth changing; in-circuit
        the difference is decisive. At add/mul granularity the saving is
        unreachable, because the circuit would be forced to trace whichever
        loop the native code happens to use.
  \item Sub-word gate parameters (BLAKE3 counters/flags): per-type pinning
        vs.\ public constant words --- decide when the BLAKE3 circuit gate
        is specified.
  \item Boolean glue tables for the recursion (byte packing,
        bit-decomposition, conditional swap): exact table inventory.
  \item Query-ladder profile for recursion-facing proofs (Slim vs.\ Fast)
        and the Merkle-caps parameter.
  \item \textbf{The jagged remainder} (named here; \emph{postponed}). At
        the end of the merged transport the verifier must evaluate each
        claim's jagged-composed weight at the final point; it cannot do so
        alone in polylog time, and the current delegation (the Frobenius
        assist) costs the verifier column-linear work $\times$ 128
        components per claim --- the dominant merged-path verify cost and
        its recursion blocker. The multipoint-twisted replacement is well
        past prototype --- a fused multipoint prover sharing the
        merged-open context, with column-parallel dispatch, landed
        2026-07-27 with measured numbers (see the capacity-free note) ---
        but the delegation remains \emph{column-linear} (active run-tree
        work improves the constant by $m/\log n_t$, not the linearity
        itself), and the soundness write-up is owed; that write-up should
        also specialize single-component (element/wiring) weights.
        Postponable for this document's milestones because:
        element claims ride the existing assist correctly (just 128$\times$
        wasteful on their components, negligible at few claims), and the
        first recursion can target the \emph{direct} path, where this
        phase does not exist.
\end{enumerate}

\end{document}
